\documentclass[oneside]{report}
\usepackage{listings,color,textcomp}
\usepackage{url}
\usepackage{graphicx}

\begin{document}

\lstset{
  language=Scheme,
  morekeywords= {
    extern, kernel, reduce, vector, void, int, u64, float, str,
    print, module, return, for, while, assert, vector-ref, length, iota, vec
  }
}

\title{Harlan User's Guide}

\maketitle

\tableofcontents

\chapter{Getting Started}

Harlan is known to build and run on the following operating systems,
but others will probably work as well.

\begin{itemize}
\item Mac OS X 10.6 (Snow Leopard)
\item Mac OS X 10.7 (Lion)
\item Gentoo Linux
\end{itemize}

Harlan requires an OpenCL implementation as well as Chez Scheme or
Petite Chez Scheme. Below are several OpenCL implementations that
should work.

\begin{itemize}
\item Intel OpenCL SDK
  \footnote{\url{http://software.intel.com/en-us/articles/opencl-sdk/}}
\item NVIDIA CUDA Toolkit
  \footnote{\url{http://developer.nvidia.com/cuda-toolkit}}
\item AMD Accelerated Parallel Processing (APP) SDK
  \footnote{\url{http://developer.amd.com/SDKS/AMDAPPSDK/Pages/default.aspx}}
\end{itemize}

Petite Cheze Scheme can be downloaded from
\url{http://www.scheme.com}.

Once all the prerequisites are installed, you can now build
Harlan. The source code for Harlan is hosted on GitHub. The following
commands will check out the source code, build Harlan, and run the
test suite.

\begin{verbatim}
    git clone git://github.iu.edu/eholk/harlan.git
    cd harlan
    make check
\end{verbatim}

If the tests are successful, you will see the following at the end of
all the output.

\begin{verbatim}
    All tests succeeded.
\end{verbatim}

The test programs are available in the \texttt{test} directory. Test
programs have the \texttt{.kfc} extension.

Make puts test binaries in the \texttt{test.bin} directory, and also
saves output from test programs here. Programs may be run directly
from this directory for easier debugging.

Harlan programs can be compiled manually as follows.

\begin{verbatim}
    ./harlanc hello.kfc
\end{verbatim}

This script has many options.  For debugging purposes, the \texttt{-v}
flag can be used to write out the intermediate results from each
compiler pass.  To see timing output, use the \texttt{-t} option.  It
is also possible to turn off verification using \texttt{-V}, as well
as turning on quiet mode with \texttt{-q}, so that no compiler
warnings or information is shown.

If \texttt{harlanc} is successful, the compiler will produce an
executable based on the input filename. This program can be run just
like any other program.

If Latex is installed, you can build the Harlan User's Guide (which
you are currently reading) as follows.

\begin{verbatim}
    make docs
\end{verbatim}


%\chapter{Using the Harlan Compiler}

\chapter{The Harlan Language}

Harlan is a language for data parallel computing that has its roots in
C and Scheme. Harlan programs resemble Scheme, but several forms have
been added to support parallel processing. One key difference is that
Harlan is statically typed. The compiler infers the type for each
expression, however, so the programmer is not burdened by type
annotations.

\section{Basics}

The traditional Hello, World! program is written as follows.

\begin{lstlisting}
(module
  (define (main)
    (print "Hello, World!")
    (return 0)))
\end{lstlisting}

This program demonstrates several characteristics of Harlan. All
programs exist within a top-level \lstinline{module} form. Harlan
modules contain a number of top-level definitions. The
\lstinline{(define (main) ...)} form defines a function called
\lstinline{main}, which is the entry point to a Harlan program.

The \lstinline{print} form is used to print any value to standard
out. In this case, the phrase ``Hello, World!'' will be printed to the
terminal. Finally, \lstinline{(return 0)} passes the value of 0 to the
calling function, which is the C convention for indicating a
successful program execution.

Programs can contain multiple procedure definitions, and procedures
can take zero or more arguments:

\begin{lstlisting}
(module
  (define (add x y)
    (return (+ x y)))

  (define (main)
    (print (add 2 3))
    (return 0)))
\end{lstlisting}

This program prints the value 5 to the terminal and then exits
successfully.

Harlan supports a number of different data types, such as integers,
floating point numbers and strings. A critical component for data
parallelism is the immutable \lstinline{vector} type. Vectors are
collections of data all sharing a single type. Vectors can contain
other vectors, but since the length of a vector is considered part of
its type, only rectangular shapes are allowed. Vectors are created as
follows.

\begin{lstlisting}
(module
  (define (main)
    (let ((x (vector 1 2)))
      (return 0))))
\end{lstlisting}

This program also demonstrates the \lstinline{let} form, which behaves
like the usual Scheme let. A vector is created and stored in the
lexical variable x, a variable that can now be used within the body of
the let.

Vectors can also be nested:

\begin{lstlisting}
(module
  (define (main)
    (let ((x (vector
              (vector 1 2)
              (vector 3 4))))
      (return 0))))
\end{lstlisting}

Harlan does not support vector manipulation, but does support
Scheme-like vector accessing with \lstinline{vector-ref}. There is
also \lstinline{length}, which returns the length of a vector. The
following program is an example of using these primitives.  

\begin{lstlisting}
(module
  (define (main)
    (let ((x (vector 0 1)))
      (for (i 0 (length x))
        (assert (= (vector-ref x i) i)))
      (return 0))))
\end{lstlisting}

The primitive \lstinline{assert} will throw an error if its condition
is violated.  It also introduces the \lstinline{for} loop. For loops
have two possible forms:

\begin{lstlisting}
  (for (<var> <start> <end>) <body>)
  (for (<var> <start> <end> <step>) <body>)
\end{lstlisting}

For loops bind the variable \lstinline{<var>} and executes
\lstinline{<body>} with the loop variable set to each value in the
range from \lstinline{<start>} up to but not including
\lstinline{<end>}, counting by \lstinline{<step>}.

These vector examples so far have not been terribly interesting, since
they do not display anything to the user. Fortunately, it is also
possible to print more complex values, including nested vectors.

\begin{lstlisting}
(module
  (define (main)
    (let ((x (vector
              (vector 1 2)
              (vector 3 4))))
      (print x)
      (return 0))))
\end{lstlisting}

This program executes and produces the following result.

\begin{verbatim}
[[1 2]
 [3 4]]
\end{verbatim}

It is also possible to use external functions.  The \lstinline{extern}
form takes the name of a procedure, as well as a type signature for
it.  The type signature is necessary because the runtime library is
written in C++ and Harlan needs to know the corresponding type.  The
available externals are described in Chapter \ref{chap:externs}.

\section{Data Parallelism}

Harlan supports data parallelism through two forms: \lstinline{kernel}
and \lstinline{reduce}.  Kernels most naturally correspond to GPU
kernel functions from CUDA or OpenCL.  The reduce operator, on the
other hand, is used to combine multiple values into a single one.

The kernel block takes a list of bindings followed by an
expression.

\begin{lstlisting}
(let ((xs (vector 1 2 3 4)))
  (kernel ((x xs))
    (+ x 1)))
\end{lstlisting}

In this example, the binding \lstinline{(x xs)} states that
\lstinline{x} ranges over the values in the vector \lstinline{xs}. For
each value of \lstinline{x}, the kernel returns the value of
\lstinline{x} plus one.  The result of the kernel expression is a
vector containing 2, 3, 4 and 5.

Kernels are allowed to refer to free variables, but cannot mutate
them.  For example, the previous example can be written like this:

\begin{lstlisting}
(let ((xs (vector 1 2 3 4)))
  (kernel ((i (iota 4)))
    (+ (vector-ref xs i) 1)))
\end{lstlisting}

Here, \lstinline{(iota 4)} expands into a vector containing 0, 1, 2
and 3. The kernel expression evaluates to the same thing as the
previous example, but the vector references are now explicit.

Multiple arguments can be given to a single kernel expression. The
kernel will run over all arguments together. For example, the
following kernel adds two vectors.

\begin{lstlisting}
(let ((xs (vector 1 2 3 4))
      (ys (vector 5 6 7 8)))
  (kernel ((x xs)
           (y ys))
    (+ x y)))
\end{lstlisting}

Kernels can be nested as well:

\begin{lstlisting}
(kernel ((a (vector 2 3)))
  (kernel ((b (vector 5 7)))
    (* a b)))  
\end{lstlisting}

This expresion computes the outer product of \lstinline{(vector 2 3)}
and \lstinline{(vector 5 7)}. If we print the result, we see we get
the expected value:

\begin{verbatim}
[[10 14]
 [15 21]]
\end{verbatim}

The other data parallel operator in Harlan is the \lstinline{reduce}
operator. This is used to combine a vector of values into a single
value. Below is a simple example.

\begin{lstlisting}
(reduce + (vector 1 1 1 1))  
\end{lstlisting}

This expression combines all the values in its vector using the
addition operator. The expression evaluates to 4.

Reduce can also be used to calculate a dot product.

\begin{lstlisting}
(reduce + (kernel ((x xs) (y ys))
            (* x y)))
\end{lstlisting}

\section{Example: The Mandelbrot Set}

Figure \ref{fig:mandelbrot-code} shows a program that calculates the
Mandelbrot set and writes it to a file.  Running this program produces
the image in Figure \ref{fig:mandelbrot}.

\begin{figure}
\begin{lstlisting}
(module
  (extern nanotime () -> u64)
  (define (main)
    (let ((start (nanotime)))
      (let ((img (kernel ((i (iota 1024)))
                   (kernel ((j (iota 1024)))
                     (let ((cr (- (/ (int->float j) 512.0) 1.5))
                           (ci (- (/ (int->float i) 512.0) 1.0))
                           (xr 0.0)
                           (xi 0.0)
                           (idx 0)
                           (escape 0))
                       (while (< idx 256)
                         (let ((xrp (+ (- (* xr xr) (* xi xi)) cr))
                               (xip (+ (* 2.0 (* xr xi)) ci)))
                           (set! xr xrp)
                           (set! xi xip)
                           (let ((m (+ (* xr xr) (* xi xi))))
                             (if (< m 4.0)
                                 (set! escape (- idx 1)))
                             (set! idx (+ idx 1)))))
                       escape)))))
        (let ((stop (nanotime)))
          (write-pgm "test.bin/mandelbrot.pgm" img)
          (print "Time to generate Mandelbrot Set in milliseconds:\n")
          (print (/ (- stop start) 1000000))
          (print "\n")
          (return 0))))))
\end{lstlisting}
  
  \caption{A program that generates the Mandelbrot Set.}
  \label{fig:mandelbrot-code}
\end{figure}

\begin{figure}
  \centering
  \includegraphics[width=\textwidth]{figures/mandelbrot.eps}
  \caption{The Mandelbrot Set, as generated by a Harlan program.}
  \label{fig:mandelbrot}
\end{figure}

\chapter{The Runtime Library}\label{chap:externs}

Harlan provides a small library of built-in functions. These are
described in this chapter. Each of these may be used in a program by
adding the appropriate \lstinline{extern} definition.

\section{nanotime}

\textbf{Prototype:} \lstinline{(extern nanotime () -> u64)}

This function returns the number of nanoseconds that has elapsed since
January 1, 1970.  It is useful for timing benchmark programs.

\section{open\_outfile and close\_outfile}

\textbf{Prototype:} \lstinline{(extern open_outfile (str) -> (ptr ofstream))}

\textbf{Prototype:} \lstinline{(extern close_outfile ((ptr ofstream)) -> void)}

Opens and closes a file that may be written to, using
\lstinline{print}.  For example,

\begin{lstlisting}
(module
  (extern open_outfile (str) -> (ptr ofstream))
  (extern close_outfile ((ptr ofstream)) -> void)
  (define (main)
    (let ((f (open_outfile "outfile.txt")))
      (print "Hello there!\n" f)
      (close_outfile f)
      (return 0))))
\end{lstlisting}

Instead of printing to the terminal, print sends its output to the open file.

\end{document}
